\documentclass[11pt,a4paper]{article}

\usepackage[margin=1in]{geometry}
\usepackage[T1]{fontenc}
\usepackage{lmodern}
\usepackage{microtype}
\usepackage{parskip}
\usepackage{xcolor}
\usepackage{listings}
\usepackage{booktabs}
\usepackage{tabularx}
\usepackage{enumitem}
\usepackage{hyperref}
\usepackage{titlesec}

\hypersetup{
  colorlinks=true,
  linkcolor=blue!60!black,
  urlcolor=blue!60!black,
  citecolor=blue!60!black
}

\definecolor{codebg}{RGB}{245,245,247}
\definecolor{codekw}{RGB}{20,90,160}
\definecolor{codestr}{RGB}{30,120,60}
\definecolor{codecom}{RGB}{120,120,120}

\lstdefinelanguage{hcl}{
  keywords={resource, variable, output, module, provider, data, locals, terraform, required_providers, required_version, true, false, null, for_each, count, depends_on},
  keywordstyle=\color{codekw}\bfseries,
  sensitive=true,
  morecomment=[l]{\#},
  morestring=[b]",
}

\lstset{
  basicstyle=\ttfamily\small,
  backgroundcolor=\color{codebg},
  frame=single,
  framerule=0pt,
  keywordstyle=\color{codekw}\bfseries,
  stringstyle=\color{codestr},
  commentstyle=\color{codecom}\itshape,
  breaklines=true,
  breakatwhitespace=true,
  showstringspaces=false,
  columns=fullflexible,
  xleftmargin=8pt,
  xrightmargin=8pt,
  aboveskip=8pt,
  belowskip=8pt,
}

\titleformat{\section}{\Large\bfseries}{\thesection}{0.7em}{}
\titleformat{\subsection}{\large\bfseries}{\thesubsection}{0.6em}{}

\title{\textbf{AIKnowledgeHub}\\[4pt]\large AWS Serverless Infrastructure Reference}
\author{AIKnowledgeHub Project}
\date{\today}

\begin{document}
\maketitle

\begin{abstract}
\noindent
AIKnowledgeHub is a serverless pipeline that scrapes fresh AI research papers from arXiv, generates plain-language summaries with Amazon Bedrock, converts those summaries to audio with Amazon Polly, stores results in DynamoDB and S3, and delivers personalised email digests to subscribers through Amazon SES. The entire back-end is defined in Terraform and runs as four AWS Lambda functions coordinated by Amazon SQS queues and Amazon EventBridge schedules. This document is a complete reference for the infrastructure: what each resource does, why it exists, how the Lambdas interact with it, and how to deploy, operate, and tear it down.
\end{abstract}

\tableofcontents
\vspace{1em}
\hrule

\section{Project Overview}

\subsection{End goal}
Deliver a personalised daily digest of new AI research papers to subscribers. Each paper includes a machine-generated summary (3--4 sentences, in plain English) and a short audio version of that summary so users can read or listen to keep up with the field passively.

\subsection{Repository layout}
The Git repository contains three top-level directories:

\begin{itemize}[leftmargin=*]
  \item \texttt{backend/} --- the original Express REST API. It still serves the front-end (\texttt{/papers}, \texttt{/subscriptions}, etc.) using MongoDB. It is unchanged by this work.
  \item \texttt{frontend/} --- the React/Vite single-page application.
  \item \texttt{infra/} --- the serverless back-end: Terraform code for all AWS resources, plus the four Lambda function sources. \textit{This document describes the contents of this folder.}
\end{itemize}

\subsection{Why serverless}
The original monolithic Express backend mixes three unrelated concerns: scraping, AI summarisation, and email. Each has very different scaling and reliability characteristics. Breaking them into four Lambdas connected by queues:
\begin{itemize}[leftmargin=*]
  \item lets each function scale independently,
  \item isolates failure (a Bedrock outage cannot stall the scraper),
  \item removes the need for always-on compute,
  \item makes retries, dead-letter queues, and backoff free.
\end{itemize}

\section{Top-Level Architecture}

\subsection{Component diagram}
\begin{lstlisting}[basicstyle=\ttfamily\footnotesize]
         +----------------------+
         |     EventBridge      |
         |   (cron schedules)   |
         +----------+-----------+
                    |
        +-----------+------------+
        |                        |
        v                        v
+-------------------+   +----------------------------+
| scraper-processor |   |   subscription-processor   |
|    (Lambda)       |   |         (Lambda)           |
+---------+---------+   +-------------+--------------+
          |                           |
     DynamoDB                    SQS email_queue
          |                           |
          v                           v
     SQS paper_queue         +---------------------+
          |                   |    email-sender     |
          v                   |  (Lambda + SES)     |
+-------------------+         +---------------------+
|   ai-processor    |
|    (Lambda)       |
+--+-----+-----+----+
   v     v     v
 Bedrock Polly  S3 (audio MP3s)
                    |
                    v
              DynamoDB (audioUrl, aiSummary)
\end{lstlisting}

\subsection{Data flow in prose}
\begin{enumerate}[leftmargin=*]
  \item EventBridge fires daily at 02:00 UTC $\rightarrow$ invokes \texttt{scraper-processor}.
  \item \texttt{scraper-processor} fetches the latest papers from arXiv for three categories (NLP, Computer Vision, Reinforcement Learning), writes each new paper to the \texttt{aikhub-papers} DynamoDB table using a conditional put (silent deduplication on \texttt{externalId}), and pushes one SQS message per topic summarising how many new papers were inserted.
  \item SQS's event-source mapping automatically invokes \texttt{ai-processor} with those messages (batches of up to five).
  \item \texttt{ai-processor} queries DynamoDB for papers in that topic that do not yet have an \texttt{audioUrl}, calls Bedrock to summarise each one, calls Polly to turn the summary into an MP3, uploads the MP3 to S3, and updates the DynamoDB item with both the summary and the audio URL.
  \item EventBridge fires daily at 08:00 UTC $\rightarrow$ invokes \texttt{subscription-processor}.
  \item \texttt{subscription-processor} scans the \texttt{aikhub-subscriptions} table, queries \texttt{aikhub-papers} by topic using the \texttt{topic-createdAt-index} GSI for each user's interests (filtered by \texttt{createdAt > lastEmailSent}), and publishes a shrunken payload of up to 20 papers to the \texttt{email\_queue}.
  \item SQS's event-source mapping invokes \texttt{email-sender}, which renders a plain-text + HTML digest and delivers it via SES.
\end{enumerate}

\section{AWS Services Used}

\begin{table}[h]
\centering
\begin{tabularx}{\linewidth}{lX}
\toprule
\textbf{Service} & \textbf{Role in AIKnowledgeHub} \\
\midrule
AWS Lambda          & Hosts the four back-end functions. \\
Amazon SQS          & Queues decouple the Lambdas; DLQs hold poison messages. \\
Amazon EventBridge  & Cron schedules for the scraper and subscription processors. \\
Amazon DynamoDB     & Persistent store for papers and subscriptions. \\
Amazon S3           & Stores generated MP3 audio summaries. \\
Amazon Bedrock      & LLM (Anthropic Claude 3 Haiku) for text summarisation. \\
Amazon Polly        & Neural TTS engine for generating MP3s. \\
Amazon SES          & Sends the digest emails. \\
Amazon VPC          & Network isolation for the Lambdas. \\
NAT Gateway         & Outbound internet access from private subnets (for arXiv). \\
VPC Endpoints       & Private paths from Lambdas to SQS, SES, and DynamoDB. \\
AWS IAM             & Single execution role granting least-privilege access. \\
Amazon CloudWatch   & Central log aggregation for every Lambda. \\
\bottomrule
\end{tabularx}
\caption{AWS services composing the infrastructure.}
\end{table}

\section{Terraform Layout}

The \texttt{infra/} folder contains one Terraform file per concern. All files are in the same module; there are no nested modules. All resources are tagged with \texttt{var.tags}.

\begin{table}[h]
\centering
\begin{tabularx}{\linewidth}{lX}
\toprule
\textbf{File} & \textbf{What it declares} \\
\midrule
\texttt{providers.tf}    & Terraform $\geq$ 1.5 and AWS provider $\sim$5.0 pinned; region from variable. \\
\texttt{variables.tf}    & Inputs: \texttt{aws\_region}, \texttt{project\_name}, \texttt{ses\_sender\_email}, \texttt{bedrock\_model\_id}, \texttt{tags}. \\
\texttt{vpc.tf}          & VPC, subnets, IGW, NAT, route tables, lambda SG. \\
\texttt{networking.tf}   & VPC interface endpoints for SQS and SES, Gateway endpoint for DynamoDB. \\
\texttt{sqs.tf}          & Two main queues and two dead-letter queues. \\
\texttt{s3.tf}           & Audio bucket with encryption, versioning, lifecycle rules. \\
\texttt{dynamodb.tf}     & \texttt{aikhub-papers} (with GSI) and \texttt{aikhub-subscriptions}. \\
\texttt{iam.tf}          & Single Lambda execution role plus inline policy. \\
\texttt{lambdas.tf}      & Four Lambda functions + SQS event-source mappings. \\
\texttt{eventbridge.tf}  & Two cron rules, two targets, two invoke permissions. \\
\texttt{monitoring.tf}   & Pre-created CloudWatch log groups. \\
\texttt{outputs.tf}      & Lambda names, queue URLs, table names, bucket name, NAT EIP. \\
\texttt{terraform.tfvars.example} & Template for secret/variable values. \\
\texttt{build.sh}        & Packages Lambda source into zips before \texttt{terraform apply}. \\
\texttt{lambda/}         & Source for all four Lambdas plus shared code. \\
\bottomrule
\end{tabularx}
\caption{Contents of the \texttt{infra/} directory.}
\end{table}

\section{Networking}

\subsection{VPC topology}
\texttt{vpc.tf} creates a fresh VPC with CIDR \texttt{10.0.0.0/16}. Two public subnets (\texttt{10.0.1.0/24}, \texttt{10.0.2.0/24}) host the NAT Gateway; two private subnets (\texttt{10.0.11.0/24}, \texttt{10.0.12.0/24}) host the VPC-attached Lambda ENIs. The subnets are spread across the first two availability zones in the chosen region.

\subsection{Routing}
The public route table routes \texttt{0.0.0.0/0} through the Internet Gateway. The private route table routes \texttt{0.0.0.0/0} through the NAT Gateway, giving Lambdas outbound internet access without making them reachable from outside.

\subsection{VPC endpoints}
Three VPC endpoints bypass the internet for common AWS calls:
\begin{itemize}[leftmargin=*]
  \item \textbf{SQS} --- Interface endpoint in both private subnets. Lambdas call SQS by private IP.
  \item \textbf{SES} --- Interface endpoint (\texttt{com.amazonaws.<region>.email}) for SES API calls.
  \item \textbf{DynamoDB} --- Gateway endpoint attached to the private route table. Free, and avoids paying for NAT bandwidth on DynamoDB traffic.
\end{itemize}

Interface endpoints create ENIs in the private subnets. These ENIs use the Lambda security group, which has a self-referencing inbound 443 rule so Lambdas can actually connect to them. (Without that rule, endpoint ENIs silently drop inbound packets.)

\subsection{Security group}
A single \texttt{aws\_security\_group.lambda} covers both Lambda ENIs and endpoint ENIs:
\begin{itemize}[leftmargin=*]
  \item \textbf{Egress:} all traffic allowed (needed for arXiv, Bedrock, Polly, SES, etc.).
  \item \textbf{Ingress:} 443/TCP from itself --- lets Lambdas reach interface endpoints.
\end{itemize}

\section{Data Model (DynamoDB)}

\subsection{\texttt{aikhub-papers}}

\begin{table}[h]
\centering
\begin{tabularx}{\linewidth}{lX}
\toprule
\textbf{Attribute} & \textbf{Meaning} \\
\midrule
\texttt{externalId} (PK, String) & arXiv paper identifier (e.g.\ \texttt{2401.12345v1}). \\
\texttt{title} (String)          & Paper title. \\
\texttt{topic} (String)          & Single primary topic: NLP, Computer Vision, or Reinforcement Learning. Used as GSI partition key. \\
\texttt{topics} (List)           & Original topic array, kept for display. \\
\texttt{summary} (String)        & Abstract from arXiv (raw). \\
\texttt{aiSummary} (String)      & LLM-generated 3--4 sentence summary (set by \texttt{ai-processor}). \\
\texttt{audioUrl} (String)       & S3 URL of the Polly-generated MP3 (set by \texttt{ai-processor}). \\
\texttt{pdfUrl} (String)         & Direct link to the paper's PDF on arXiv. \\
\texttt{source} (String)         & \texttt{"arXiv"}. \\
\texttt{publishedAt} (String)    & ISO-8601 original publish date. \\
\texttt{createdAt} (String)      & ISO-8601 insertion time. GSI sort key. \\
\texttt{views} (Number)          & View count (updated by the front-end path; 0 at insertion). \\
\bottomrule
\end{tabularx}
\caption{Attributes of the \texttt{aikhub-papers} table.}
\end{table}

\paragraph{Global secondary index.} \texttt{topic-createdAt-index} has partition key \texttt{topic} and sort key \texttt{createdAt}, projecting all attributes. This supports the two hot queries:
\begin{enumerate}[leftmargin=*]
  \item \texttt{ai-processor} reads unprocessed papers for a topic.
  \item \texttt{subscription-processor} reads papers for a topic newer than a user's \texttt{lastEmailSent}.
\end{enumerate}

\paragraph{Billing.} Both tables use \texttt{PAY\_PER\_REQUEST}: no capacity planning, costs scale with traffic, ideal for bursty / low-volume use.

\subsection{\texttt{aikhub-subscriptions}}

\begin{table}[h]
\centering
\begin{tabularx}{\linewidth}{lX}
\toprule
\textbf{Attribute} & \textbf{Meaning} \\
\midrule
\texttt{email} (PK, String)   & Subscriber email address. \\
\texttt{subscribedTopics} (List) & Array of topic names the user is interested in. \\
\texttt{lastEmailSent} (String) & ISO-8601 of the most recent digest delivery. Used to avoid re-sending. \\
\bottomrule
\end{tabularx}
\caption{Attributes of the \texttt{aikhub-subscriptions} table.}
\end{table}

\section{Queue Topology (SQS)}

\begin{itemize}[leftmargin=*]
  \item \textbf{\texttt{aikhub-paper-processing}} --- receives scraper fan-out messages; consumed by \texttt{ai-processor}. Visibility timeout 300\,s (matches the Lambda's max runtime).
  \item \textbf{\texttt{aikhub-email-queue}} --- receives one message per subscriber-digest; consumed by \texttt{email-sender}. Visibility timeout 60\,s.
  \item Each has a matching \textbf{dead-letter queue} (\texttt{*-dlq}) with 14-day retention and a redrive policy of \texttt{maxReceiveCount = 3}.
\end{itemize}

Both main queues are wired to their consumer via \texttt{aws\_lambda\_event\_source\_mapping} with a \texttt{batch\_size} of 5. Failed items can be reported back through the standard \texttt{batchItemFailures} response; this is used by \texttt{email-sender} so that a failure for one recipient does not reprocess the entire batch.

\section{Scheduling (EventBridge)}

Two \texttt{aws\_cloudwatch\_event\_rule} resources:

\begin{center}
\begin{tabular}{lll}
\toprule
\textbf{Name} & \textbf{Cron} & \textbf{Target} \\
\midrule
\texttt{aikhub-daily-scrape} & \texttt{cron(0 2 * * ? *)} & \texttt{scraper-processor} \\
\texttt{aikhub-daily-email}  & \texttt{cron(0 8 * * ? *)} & \texttt{subscription-processor} \\
\bottomrule
\end{tabular}
\end{center}

Each rule has a corresponding \texttt{aws\_lambda\_permission} granting \texttt{events.amazonaws.com} invoke rights for just that rule's ARN.

\section{IAM}

A single execution role \texttt{aikhub-lambda-exec} is assumed by all four Lambdas. Its inline policy grants only the actions each Lambda actually uses:

\begin{itemize}[leftmargin=*]
  \item \texttt{logs:CreateLogGroup, CreateLogStream, PutLogEvents} --- CloudWatch Logs.
  \item \texttt{sqs:SendMessage, ReceiveMessage, DeleteMessage, GetQueueAttributes} --- queue I/O.
  \item \texttt{ses:SendEmail, SendRawEmail} --- email delivery.
  \item \texttt{bedrock:InvokeModel} --- LLM summarisation.
  \item \texttt{polly:SynthesizeSpeech} --- text-to-speech.
  \item \texttt{s3:PutObject, GetObject} --- scoped to \texttt{\$\{audio\_bucket.arn\}/*}.
  \item \texttt{ec2:CreateNetworkInterface, DescribeNetworkInterfaces, DeleteNetworkInterface} --- required for VPC-attached Lambdas.
  \item \texttt{dynamodb:PutItem, GetItem, UpdateItem, DeleteItem, Query, Scan, BatchGetItem, BatchWriteItem} --- scoped to both tables' ARNs and \texttt{papers/index/*} for GSI access.
\end{itemize}

Bedrock and Polly actions are intentionally scoped to \texttt{"*"} because service-level ARNs are clumsy and these APIs do not support fine-grained resource policies usefully. The other high-risk actions (S3, DynamoDB) are ARN-scoped.

\section{Storage (S3)}

The \texttt{aikhub-audio-summaries} bucket holds all Polly-generated MP3s. Hardened with:
\begin{itemize}[leftmargin=*]
  \item \textbf{Public access block:} all four flags on.
  \item \textbf{Server-side encryption:} AES-256 default.
  \item \textbf{Versioning:} enabled.
  \item \textbf{Lifecycle rule:} non-current versions expire after 30 days; stalled multipart uploads abort after 7.
\end{itemize}

Object keys use the sanitised paper ID: \texttt{audio/\{externalId\}.mp3}. Because the bucket is private, any front-end that wants to play these files needs either a pre-signed URL or a CloudFront distribution with an Origin Access Control. Wiring that up is out of scope for the current infrastructure.

\section{Lambda Functions}

\subsection{\texttt{scraper-processor}}

\textbf{Trigger:} EventBridge cron (daily at 02:00 UTC) or manual \texttt{aws lambda invoke}. \\
\textbf{Runtime:} Node.js 20, 256\,MB, 120\,s timeout, VPC-attached. \\
\textbf{Environment:} \texttt{PAPERS\_TABLE}, \texttt{PAPER\_QUEUE\_URL}. \\
\textbf{Code:} \texttt{lambda/scraper-processor/index.js} and \texttt{services/scraper.service.js}.

\paragraph{What it does.}
\begin{enumerate}[leftmargin=*]
  \item For each of three arXiv categories (\texttt{cs.CL}, \texttt{cs.CV}, \texttt{cs.LG}), fetches up to 10 newest papers via the public arXiv API.
  \item Parses the Atom XML response with \texttt{fast-xml-parser} and extracts title, summary, PDF link, and external ID.
  \item Calls \texttt{PutCommand} on DynamoDB with \texttt{ConditionExpression: "attribute\_not\_exists(externalId)"}. Duplicates throw \texttt{ConditionalCheckFailedException} which is caught and counted as \texttt{skipped} --- no errors, no cost for rewriting existing rows.
  \item For every category that inserted at least one new paper, pushes a JSON message to \texttt{paper\_queue}: \texttt{\{ action, topic, category, insertedCount \}}. The downstream consumer uses \texttt{topic} to know which papers to process.
\end{enumerate}

\subsection{\texttt{ai-processor}}

\textbf{Trigger:} SQS event-source mapping on \texttt{paper\_queue}, batch size 5. \\
\textbf{Runtime:} Node.js 20, 512\,MB, 300\,s timeout, VPC-attached. \\
\textbf{Environment:} \texttt{PAPERS\_TABLE}, \texttt{S3\_BUCKET}, \texttt{BEDROCK\_MODEL\_ID}.

\paragraph{What it does.} For each SQS record:
\begin{enumerate}[leftmargin=*]
  \item Queries \texttt{topic-createdAt-index} for the named topic, with \texttt{FilterExpression: "attribute\_not\_exists(audioUrl)"}. Only papers that still need processing come back. A small \texttt{Limit} multiplier compensates for the filter happening post-query.
  \item For each returned paper (double-check \texttt{audioUrl} absent for idempotency on SQS retry):
    \begin{itemize}[leftmargin=*]
      \item \textbf{Bedrock.} Calls \texttt{ConverseCommand} with the configured inference profile (default \texttt{eu.anthropic.claude-3-haiku-20240307-v1:0}). Converse normalises request format across Anthropic, Amazon Nova, Llama, etc., so the model can be swapped in \texttt{terraform.tfvars} without code changes.
      \item \textbf{Polly.} \texttt{SynthesizeSpeech} with neural voice \texttt{Joanna}, MP3 output.
      \item \textbf{S3.} \texttt{PutObject} under \texttt{audio/\{sanitizedId\}.mp3}.
      \item \textbf{DynamoDB.} \texttt{UpdateItem} on the paper, setting \texttt{aiSummary} and \texttt{audioUrl}.
    \end{itemize}
\end{enumerate}

If any step throws, the whole batch item fails and SQS re-delivers it up to three times before routing to \texttt{paper\_dlq}. Because the \texttt{audioUrl} check is re-tested per paper, a retry after a partial batch will not re-summarise or re-upload already-done papers.

\subsection{\texttt{subscription-processor}}

\textbf{Trigger:} EventBridge cron (daily at 08:00 UTC). \\
\textbf{Runtime:} Node.js 20, 256\,MB, 120\,s timeout, VPC-attached. \\
\textbf{Environment:} \texttt{PAPERS\_TABLE}, \texttt{SUBSCRIPTIONS\_TABLE}, \texttt{EMAIL\_QUEUE\_URL}.

\paragraph{What it does.}
\begin{enumerate}[leftmargin=*]
  \item \texttt{Scan}s the \texttt{aikhub-subscriptions} table. Scan is acceptable here because the subscriber table is small and only iterated once per day.
  \item For every user, queries \texttt{topic-createdAt-index} in parallel for each of their \texttt{subscribedTopics}. The key condition is \texttt{topic = :t AND createdAt > :lastEmailSent} when the user already received an email, otherwise \texttt{topic = :t}.
  \item If the user has no new papers, falls back to the top-5 most-viewed papers in their topics so they always get something.
  \item Trims the result to 20 papers and 500-character summaries to stay safely under the SQS message size limit of 256\,KB. Each slimmed paper item carries only \texttt{externalId}, \texttt{title}, \texttt{topic}, \texttt{summary} (preferring \texttt{aiSummary}), \texttt{pdfUrl}.
  \item Sends the payload to \texttt{email\_queue}.
  \item \texttt{UpdateItem}s the user's \texttt{lastEmailSent} to the current ISO timestamp.
\end{enumerate}

\subsection{\texttt{email-sender}}

\textbf{Trigger:} SQS event-source mapping on \texttt{email\_queue}, batch size 5. \\
\textbf{Runtime:} Node.js 20, 128\,MB, 30\,s timeout, \emph{not} VPC-attached (SES public endpoint is fine; avoids ENI overhead). \\
\textbf{Environment:} \texttt{SES\_SENDER\_EMAIL}.

\paragraph{What it does.}
\begin{enumerate}[leftmargin=*]
  \item Renders a plain-text and HTML body from the slimmed paper list.
  \item Calls \texttt{SESClient.SendEmail}.
  \item Per record, appends \texttt{\{ messageId, status \}} to a results buffer.
  \item Returns a \texttt{batchItemFailures} array so that partial batch failures only retry the failed items, not the whole batch.
\end{enumerate}

\subsection{Shared code}
\texttt{lambda/shared/dynamo.js} exports a pre-configured DynamoDB Document Client:
\begin{lstlisting}[language=java]
import { DynamoDBClient } from "@aws-sdk/client-dynamodb";
import { DynamoDBDocumentClient } from "@aws-sdk/lib-dynamodb";

const baseClient = new DynamoDBClient({});
export const ddb = DynamoDBDocumentClient.from(baseClient, {
  marshallOptions: {
    removeUndefinedValues: true,
    convertEmptyValues: false,
  },
});
\end{lstlisting}
No caching is needed --- the DynamoDB client is stateless, unlike the previous Mongoose connection. This file is copied into every Lambda that needs it by \texttt{build.sh}.

\section{Packaging: \texttt{build.sh}}

Terraform uploads a Lambda as a zip. \texttt{build.sh} produces the zips each deploy. For every Lambda:

\begin{enumerate}[leftmargin=*]
  \item Wipes and recreates \texttt{infra/builds/}.
  \item Copies the Lambda source into a staging directory.
  \item If the Lambda imports from \texttt{./shared}, copies \texttt{lambda/shared/} into the staging folder.
  \item Runs \texttt{npm install --omit=dev --no-audit --no-fund --silent} in the staging folder to pull in production dependencies into \texttt{node\_modules/}.
  \item Zips the staging contents into \texttt{builds/<lambda-name>.zip}. On Linux/macOS this uses \texttt{zip}; on Git Bash (Windows) it falls back to PowerShell's \texttt{Compress-Archive}, with paths converted through \texttt{cygpath -w}.
  \item Deletes the staging folder.
\end{enumerate}

Terraform detects source changes via \texttt{source\_code\_hash = filebase64sha256(...)}. The next \texttt{terraform apply} automatically uploads changed zips.

\section{Deployment Workflow}

\subsection{Prerequisites}
\begin{itemize}[leftmargin=*]
  \item Terraform $\geq$ 1.5, Node.js $\geq$ 18, AWS CLI v2.
  \item AWS credentials configured locally (\texttt{aws configure}; \emph{never} pasted into repo files).
  \item SES sender identity verified in the target region.
  \item Bedrock model access: for Anthropic Claude models, the use-case form must be submitted once per account via the Bedrock console.
\end{itemize}

\subsection{Config}
Copy \texttt{terraform.tfvars.example} to \texttt{terraform.tfvars} and fill:
\begin{lstlisting}[language=hcl]
aws_region       = "eu-west-1"
project_name     = "aikhub"
ses_sender_email = "you@yourdomain.com"
bedrock_model_id = "eu.anthropic.claude-3-haiku-20240307-v1:0"
\end{lstlisting}

\texttt{terraform.tfvars} is \texttt{.gitignore}d. The \texttt{.example} file is tracked and must contain only placeholder values.

\subsection{Apply}
\begin{lstlisting}[basicstyle=\ttfamily\small]
cd infra
./build.sh
terraform init
terraform apply
\end{lstlisting}

The first apply provisions ${\sim}45$ resources in 3--5 minutes. The NAT Gateway is the slowest; VPC endpoints take ${\sim}60$ seconds each.

\subsection{Teardown}
\begin{lstlisting}[basicstyle=\ttfamily\small]
terraform destroy
\end{lstlisting}

The NAT Gateway bills hourly, so teardown matters. \emph{Warning:} VPC-attached Lambdas produce Hyperplane ENIs that can take 20--45 minutes to release, blocking security-group destruction. Plan teardowns with that in mind.

\section{Observability}

Each Lambda has a pre-created \texttt{aws\_cloudwatch\_log\_group} at \texttt{/aws/lambda/aikhub-<name>} with 14-day retention. Terraform owning the log group (rather than letting Lambda auto-create it on first invocation) enforces retention and avoids the classic drift where the log group exists but is untracked.

Live tailing (Git Bash on Windows needs the \texttt{MSYS\_NO\_PATHCONV} escape to prevent path mangling):

\begin{lstlisting}[basicstyle=\ttfamily\small]
MSYS_NO_PATHCONV=1 aws logs tail /aws/lambda/aikhub-ai-processor \
  --region eu-west-1 --follow
\end{lstlisting}

\section{Operational Scenarios}

\subsection{Manual invocation / smoke test}
\begin{lstlisting}[basicstyle=\ttfamily\small]
aws lambda invoke --function-name aikhub-scraper-processor \
  --region eu-west-1 /tmp/out.json
cat /tmp/out.json
\end{lstlisting}

\subsection{Triggering the pipeline with a seeded SQS message}
Useful when the scraper has nothing new to insert (so no natural SQS fan-out):
\begin{lstlisting}[basicstyle=\ttfamily\small]
aws sqs send-message --region eu-west-1 \
  --queue-url <paper_queue_url> \
  --message-body '{"action":"process_new_papers","topic":"NLP","insertedCount":2}'
\end{lstlisting}

\subsection{Inspecting DynamoDB}
\begin{lstlisting}[basicstyle=\ttfamily\small]
aws dynamodb scan --table-name aikhub-papers --region eu-west-1 --select COUNT
aws dynamodb scan --table-name aikhub-papers --region eu-west-1 \
  --filter-expression "attribute_exists(audioUrl)" --select COUNT
\end{lstlisting}

\subsection{Seeding a test subscriber}
\begin{lstlisting}[basicstyle=\ttfamily\small]
aws dynamodb put-item --table-name aikhub-subscriptions --region eu-west-1 \
  --item '{"email":{"S":"you@example.com"},
           "subscribedTopics":{"L":[{"S":"NLP"}]}}'
\end{lstlisting}

Remember: SES accounts start in the \emph{sandbox}, where the recipient address must also be verified. Request production access via the SES dashboard to send to arbitrary recipients.

\section{Test Execution and Results}

The following tests were executed against a live deployment in \texttt{eu-west-1} on 2026-04-15/16. All AWS resources were provisioned via \texttt{terraform apply} (44 resources initial deploy, 3 added + 5 changed after the DynamoDB migration).

\subsection{Test 1: Scraper --- arXiv ingestion}

\textbf{Method:} Invoked the scraper Lambda directly.
\begin{lstlisting}[basicstyle=\ttfamily\small]
aws lambda invoke --function-name aikhub-scraper-processor \
  --region eu-west-1 --payload '{}' \
  --cli-binary-format raw-in-base64-out /tmp/scraper.json
\end{lstlisting}

\textbf{Result:} \texttt{StatusCode: 200}. DynamoDB scan confirmed \textbf{27 papers} inserted into \texttt{aikhub-papers} table. Each paper had \texttt{externalId}, \texttt{title}, \texttt{summary}, \texttt{topic}, \texttt{topics}, \texttt{pdfUrl}, and \texttt{createdAt} populated. SQS message sent to \texttt{aikhub-paper-processing} queue triggering the ai-processor.

\textbf{Verdict:} \textcolor{green!60!black}{\textbf{PASS}}

\subsection{Test 2: AI Processor --- summary generation + audio}

\textbf{Method:} Bedrock daily token quota was exhausted on the new AWS account, so a local Ollama (llama3) script was used as a drop-in replacement. The script mirrors the Lambda's exact behaviour: query DynamoDB for unprocessed papers, generate summary, synthesize audio via Polly, upload MP3 to S3, update DynamoDB.

\begin{lstlisting}[basicstyle=\ttfamily\small]
MAX_PAPERS=3 TOPIC=NLP node scripts/local-ai-processor.mjs
\end{lstlisting}

\textbf{Result:}
\begin{itemize}[leftmargin=*]
  \item 3 papers processed (\texttt{2604.13018v1}, \texttt{2604.13016v1}, \texttt{2604.12911v1})
  \item Summaries: 637, 661, 765 characters respectively
  \item Audio files: 248\,KB, 235\,KB, 265\,KB (Polly neural, Joanna voice)
  \item S3 upload confirmed: \texttt{s3://aikhub-audio-summaries/audio/<id>.mp3}
  \item DynamoDB updated: \texttt{aiSummary} and \texttt{audioUrl} fields populated
\end{itemize}

\textbf{Verdict:} \textcolor{green!60!black}{\textbf{PASS}} (Polly, S3, and DynamoDB used real AWS; only the LLM call was local)

\subsection{Test 3: Subscription Processor --- fan-out to email queue}

\textbf{Method:} Seeded a test subscriber, then invoked the subscription Lambda.
\begin{lstlisting}[basicstyle=\ttfamily\small]
aws dynamodb put-item --table-name aikhub-subscriptions \
  --region eu-west-1 \
  --item '{"email":{"S":"muhammad.kobeissi@gmail.com"},
           "subscribedTopics":{"L":[{"S":"NLP"}]}}'

aws lambda invoke --function-name aikhub-subscription-processor \
  --region eu-west-1 --payload '{}' \
  --cli-binary-format raw-in-base64-out /tmp/out.json
\end{lstlisting}

\textbf{Result:} \texttt{StatusCode: 200}, \texttt{\{"success":true\}}. CloudWatch logs showed 859\,ms execution, no errors. Message enqueued to \texttt{aikhub-email-processing} for the email-sender.

\textbf{Verdict:} \textcolor{green!60!black}{\textbf{PASS}}

\subsection{Test 4: Email Sender --- SES delivery}

\textbf{Method:} Triggered automatically by the SQS message from Test~3.

\textbf{Initial result:} \texttt{ERROR: p.topic.join is not a function}. Root cause: scraper writes \texttt{topic} as a string, but the email template called \texttt{p.topic.join(", ")} expecting an array.

\textbf{Fix applied:} Added a \texttt{topicList()} helper that handles both string and array forms. Rebuilt and redeployed.

\textbf{Re-test result:} \texttt{StatusCode: 200}, duration 1056\,ms, no errors. Email received in Gmail inbox with subject ``\texttt{AIKnowledgeHub: 3 new papers in NLP}''. HTML body contained paper titles, topics, summaries, and PDF links.

\textbf{Verdict:} \textcolor{green!60!black}{\textbf{PASS}} (after bug fix)

\subsection{Test Summary}

\begin{tabularx}{\textwidth}{lccX}
\toprule
\textbf{Stage} & \textbf{Lambda} & \textbf{Result} & \textbf{Notes} \\
\midrule
Scraper          & scraper-processor       & PASS & 27 papers ingested from arXiv \\
AI Processor     & ai-processor            & PASS & Tested with local Ollama (Bedrock quota blocked) \\
Subscription     & subscription-processor  & PASS & Fan-out to SQS confirmed \\
Email Sender     & email-sender            & PASS & Bug fixed, email delivered via SES \\
\bottomrule
\end{tabularx}

\section{Teammate Setup Guide: Testing the \texttt{aws-terraform-setup} Branch}

This section provides a complete step-by-step guide for a teammate who wants to clone the branch and test the full infrastructure from scratch.

\subsection{Prerequisites}

\begin{enumerate}[leftmargin=*]
  \item \textbf{AWS CLI v2} installed and configured with credentials that have admin-level permissions:
    \begin{lstlisting}[basicstyle=\ttfamily\small]
aws configure
# Enter: Access Key ID, Secret Access Key, region: eu-west-1, output: json
    \end{lstlisting}
  \item \textbf{Terraform} $\geq$ 1.5 installed: \url{https://developer.hashicorp.com/terraform/install}
  \item \textbf{Node.js} $\geq$ 18.x (for \texttt{npm install} during Lambda packaging)
  \item \textbf{Git Bash} (on Windows) or a Unix shell (macOS/Linux)
  \item \textbf{zip} utility. On Git Bash: \texttt{pacman -S zip} or verify \texttt{which zip} returns a path. Without it, builds fall back to PowerShell \texttt{Compress-Archive} (much slower).
\end{enumerate}

\subsection{Step 1: Clone and switch to the branch}

\begin{lstlisting}[basicstyle=\ttfamily\small]
git clone https://github.com/JadR05/AIKnowledgeHub.git
cd AIKnowledgeHub/infra
git checkout aws-terraform-setup
\end{lstlisting}

\subsection{Step 2: Configure Terraform variables}

\begin{lstlisting}[basicstyle=\ttfamily\small]
cp terraform.tfvars.example terraform.tfvars
\end{lstlisting}

Edit \texttt{terraform.tfvars} and set:
\begin{itemize}[leftmargin=*]
  \item \texttt{ses\_sender\_email} --- an email address you own (you will verify it in Step~4)
  \item \texttt{aws\_region} --- keep \texttt{"eu-west-1"} or change to your preferred region
  \item \texttt{bedrock\_model\_id} --- keep the default or change (see \S{}\ref{sec:swapping-models} below)
\end{itemize}

\subsection{Step 3: Verify SES sender identity}

\begin{enumerate}[leftmargin=*]
  \item Go to the \textbf{AWS Console} $\to$ \textbf{SES} $\to$ \textbf{Verified Identities}
  \item Click \textbf{Create Identity} $\to$ choose \textbf{Email address}
  \item Enter the same email from \texttt{ses\_sender\_email}
  \item Open your inbox and click the verification link from AWS
\end{enumerate}

Note: SES starts in \emph{sandbox mode}. Both sender and recipient must be verified. Since the test emails go from your address to itself, verifying one address is sufficient.

\subsection{Step 4: Enable Bedrock model access}

\begin{enumerate}[leftmargin=*]
  \item AWS Console $\to$ \textbf{Amazon Bedrock} $\to$ \textbf{Model access}
  \item Request access to \textbf{Anthropic Claude 3 Haiku} (or whichever model you configured)
  \item Fill in the use-case form when prompted (approval is usually instant for Haiku)
  \item New accounts may have low daily token quotas --- request an increase via \textbf{Service Quotas} if needed
\end{enumerate}

\subsection{Step 5: Build Lambda packages}

\begin{lstlisting}[basicstyle=\ttfamily\small]
chmod +x build.sh
./build.sh
\end{lstlisting}

This creates four zip files in \texttt{builds/}:
\begin{itemize}[leftmargin=*]
  \item \texttt{scraper-processor.zip}
  \item \texttt{ai-processor.zip}
  \item \texttt{subscription-processor.zip}
  \item \texttt{email-sender.zip}
\end{itemize}

\subsection{Step 6: Deploy infrastructure}

\begin{lstlisting}[basicstyle=\ttfamily\small]
terraform init
terraform plan        # review what will be created
terraform apply       # type "yes" to confirm
\end{lstlisting}

This provisions ${\approx}$44 AWS resources (VPC, subnets, NAT, endpoints, DynamoDB tables, SQS queues, Lambdas, IAM roles, S3 bucket, EventBridge schedules, CloudWatch log groups).

Terraform outputs the key resource names/URLs at the end. Note the \texttt{paper\_queue\_url} and \texttt{email\_queue\_url} for manual testing.

\subsection{Step 7: Test the scraper}

\begin{lstlisting}[basicstyle=\ttfamily\small]
aws lambda invoke --function-name aikhub-scraper-processor \
  --region eu-west-1 --payload '{}' \
  --cli-binary-format raw-in-base64-out /tmp/scraper.json
cat /tmp/scraper.json
\end{lstlisting}

Verify papers landed in DynamoDB:
\begin{lstlisting}[basicstyle=\ttfamily\small]
aws dynamodb scan --table-name aikhub-papers \
  --region eu-west-1 --select COUNT
\end{lstlisting}

Expected: 20--30 papers.

\subsection{Step 8: Test AI processing}

The scraper automatically sends an SQS message that triggers \texttt{ai-processor}. Wait 1--2 minutes, then check:
\begin{lstlisting}[basicstyle=\ttfamily\small]
aws dynamodb scan --table-name aikhub-papers \
  --region eu-west-1 \
  --filter-expression "attribute_exists(audioUrl)" \
  --select COUNT
\end{lstlisting}

If the count is 0 after several minutes, check logs:
\begin{lstlisting}[basicstyle=\ttfamily\small]
MSYS_NO_PATHCONV=1 aws logs tail \
  /aws/lambda/aikhub-ai-processor --region eu-west-1 --since 10m
\end{lstlisting}

Common issue: Bedrock quota exhausted on new accounts. Wait for UTC midnight reset or request a quota increase.

\subsection{Step 9: Test subscription + email delivery}

Seed a test subscriber:
\begin{lstlisting}[basicstyle=\ttfamily\small]
aws dynamodb put-item --table-name aikhub-subscriptions \
  --region eu-west-1 \
  --item '{"email":{"S":"YOUR_VERIFIED_EMAIL"},
           "subscribedTopics":{"L":[{"S":"NLP"}]}}'
\end{lstlisting}

Trigger the subscription processor:
\begin{lstlisting}[basicstyle=\ttfamily\small]
aws lambda invoke \
  --function-name aikhub-subscription-processor \
  --region eu-west-1 --payload '{}' \
  --cli-binary-format raw-in-base64-out /tmp/sub.json
\end{lstlisting}

Check email-sender logs:
\begin{lstlisting}[basicstyle=\ttfamily\small]
MSYS_NO_PATHCONV=1 aws logs tail \
  /aws/lambda/aikhub-email-sender --region eu-west-1 --since 5m
\end{lstlisting}

Check your inbox (and spam folder) for the digest email.

\subsection{Step 10: Verify S3 audio files}

\begin{lstlisting}[basicstyle=\ttfamily\small]
aws s3 ls s3://aikhub-audio-summaries/audio/ --region eu-west-1
\end{lstlisting}

To listen to a file (generates a temporary URL valid for 1 hour):
\begin{lstlisting}[basicstyle=\ttfamily\small]
aws s3 presign s3://aikhub-audio-summaries/audio/<PAPER_ID>.mp3 \
  --region eu-west-1 --expires-in 3600
\end{lstlisting}

\subsection{Step 11: Tear down (when done testing)}

\begin{lstlisting}[basicstyle=\ttfamily\small]
terraform destroy    # type "yes" to confirm
\end{lstlisting}

\textbf{Important:} The NAT Gateway costs ${\approx}$\$1.08/day even when idle. Always destroy the stack when you are not actively testing, unless you intend to keep it running.

Note: Security group destruction can take 20--45 minutes due to Lambda Hyperplane ENI release delays. This is normal --- do not cancel the destroy.

\subsection{Swapping Bedrock Models}
\label{sec:swapping-models}

The Converse API is model-agnostic. Change the model in \texttt{terraform.tfvars}:
\begin{lstlisting}[language=hcl]
# Anthropic Claude 3 Haiku (default)
bedrock_model_id = "eu.anthropic.claude-3-haiku-20240307-v1:0"
# Amazon Nova Micro (cheapest)
bedrock_model_id = "eu.amazon.nova-micro-v1:0"
\end{lstlisting}
Run \texttt{terraform apply} --- no rebuild needed, just an environment variable update.

\subsection{Troubleshooting}

\begin{tabularx}{\textwidth}{lX}
\toprule
\textbf{Symptom} & \textbf{Fix} \\
\midrule
\texttt{TimeoutError ETIMEDOUT} on Lambda & VPC endpoint or security group misconfigured. Verify SG has self-referencing ingress on port 443. \\
\texttt{ThrottlingException} from Bedrock & Daily token quota exhausted. Wait for UTC midnight reset or request increase via Service Quotas. \\
\texttt{ResourceNotFoundException} from Bedrock & Model access not granted. Complete the use-case form in Bedrock console. \\
\texttt{MessageRejected} from SES & Recipient not verified (sandbox mode). Verify the recipient email or request production access. \\
Build hangs on \texttt{Compress-Archive} & Install native \texttt{zip}: \texttt{pacman -S zip} in Git Bash. \\
\texttt{/aws/lambda/...} path mangled on Git Bash & Prefix command with \texttt{MSYS\_NO\_PATHCONV=1}. \\
SG destroy takes 20+ minutes & Normal. Lambda Hyperplane ENIs take 20--45 min to release. Do not cancel. \\
\bottomrule
\end{tabularx}

\section{Known Limitations and Future Work}

\begin{itemize}[leftmargin=*]
  \item \textbf{No public HTTP API in Terraform.} The front-end still talks to the Express back-end for paper reads and subscription writes. An \texttt{aws\_apigatewayv2\_api} + a small read/write Lambda against DynamoDB would close this gap.
  \item \textbf{S3 audio links are private.} Emails reference \texttt{https://<bucket>.s3.amazonaws.com/audio/...}, which 403s because public access is blocked. Either switch to pre-signed URLs with a sensible TTL or front the bucket with CloudFront + OAC.
  \item \textbf{No unsubscribe flow.} Production email delivery needs a one-click unsubscribe link and a token scheme, plus double opt-in for new subscriptions.
  \item \textbf{No CI/CD.} A GitHub Actions workflow could run \texttt{./build.sh} plus \texttt{terraform plan} on pull requests.
  \item \textbf{No CloudWatch alarms.} Recommended: DLQ depth, Lambda error rate, Bedrock throttle rate.
  \item \textbf{Bedrock IAM scope.} \texttt{bedrock:InvokeModel} is currently \texttt{Resource = "*"}. It could be narrowed to the configured inference profile ARN.
  \item \textbf{New-account Bedrock quotas.} New AWS accounts in eu-west-1 have a low daily token cap. First-time testing can easily exhaust it; request a quota increase via Service Quotas before onboarding real traffic.
\end{itemize}

\section{Cost Snapshot}

Back-of-envelope for a test that processes $N$ papers per day:

\begin{itemize}[leftmargin=*]
  \item \textbf{Lambda:} Free tier covers 1\,M invocations/month. This workload is well under that.
  \item \textbf{DynamoDB on-demand:} \$1.25 per million write requests, \$0.25 per million reads. Trivial at this volume.
  \item \textbf{SQS:} First 1\,M requests/month are free.
  \item \textbf{S3:} \$0.023 per GB-month. One MP3 is ${\approx}$200\,KB; 30 papers/day = ${\approx}$6\,MB/day.
  \item \textbf{Bedrock (Claude 3 Haiku):} ${\approx}$\$0.0006 per paper.
  \item \textbf{Polly (neural):} \$16 per 1\,M characters, so ${\approx}$\$0.005 per summary.
  \item \textbf{NAT Gateway:} Flat ${\approx}$\$0.048/hour in eu-west-1, so ${\approx}$\$35/month whether idle or busy. This dominates the running cost; destroy the stack when not testing.
  \item \textbf{VPC interface endpoints:} ${\approx}$\$0.011/hour each, two of them ${\approx}$\$16/month. The DynamoDB Gateway endpoint is free.
\end{itemize}

The fixed VPC costs (NAT + interface endpoints) make this stack expensive to leave idle. For pure testing, a NAT-less ``Mode A'' variant (no VPC, Lambdas on the public AWS network, DynamoDB over public endpoints) can be created by removing \texttt{vpc\_config} from the Lambdas and deleting \texttt{vpc.tf} / \texttt{networking.tf}; cost drops to effectively zero when idle.

\section{Post-Launch Updates}

The sections above describe the initial deployment. Between the first
\texttt{terraform apply} and the production demo, several additions and
refactors were made. This section captures them. Where it disagrees with
an earlier chapter, this one is authoritative.

\subsection{New: Public compute layer (EC2 + ALB + ASG + ECR)}

The front-end was originally served from the old Express host; for the
production deploy a fully managed compute tier was added and codified in
Terraform.

\begin{itemize}[leftmargin=*]
  \item \textbf{\texttt{ecr.tf}} --- two Elastic Container Registry
        repositories (\texttt{aikhub-backend}, \texttt{aikhub-frontend})
        each with a lifecycle policy that retains the last five images.
  \item \textbf{\texttt{ec2.tf}} --- Launch Template (Amazon Linux 2023,
        \texttt{t3.micro}, public-IP disabled), an Auto Scaling Group
        (\texttt{desired=2}, \texttt{min=1}, \texttt{max=4}) spanning the
        two private subnets, an internet-facing Application Load
        Balancer, a target group with ELB health checks, and a
        \texttt{forward} listener on HTTP:80.
  \item User-data on each instance installs Docker, logs in to ECR via
        the instance profile, and pulls the backend + frontend images.
        Nginx in the frontend container reverse-proxies API traffic to
        the backend container.
  \item An ASG \emph{instance refresh} is configured (rolling, min 50\%
        healthy) so that pushing a new image tag and starting a refresh
        is enough to roll out updated code without manual intervention.
  \item The EC2 IAM role gets both \texttt{AmazonEC2ContainerRegistryReadOnly}
        and \texttt{AmazonSSMManagedInstanceCore} attached --- the latter
        lets operators shell into private-subnet instances through
        Systems Manager Session Manager, avoiding bastion hosts.
\end{itemize}

\subsection{New: HTTPS via CloudFront (\texttt{cloudfront.tf})}

The ALB only listens on HTTP:80; terminating TLS directly on the ALB
requires a cert for a hostname the account controls. To get browser-trusted
HTTPS without buying a domain, a CloudFront distribution was placed in
front of the ALB:

\begin{itemize}[leftmargin=*]
  \item Origin: \texttt{aws\_lb.app.dns\_name}, \texttt{origin\_protocol\_policy = "http-only"}.
  \item Default behaviour: \texttt{redirect-to-https}, all HTTP methods
        forwarded, compression enabled.
  \item AWS managed policies: \texttt{CachingDisabled} (dynamic API, no
        caching) and \texttt{AllViewerExceptHostHeader} (forward
        everything except \texttt{Host} to avoid confusing the ALB's
        host-based routing).
  \item \texttt{cloudfront\_default\_certificate = true} --- the
        distribution inherits the \texttt{*.cloudfront.net} cert, so no
        ACM work and no domain purchase are required.
  \item Output: \texttt{cloudfront\_url} --- the canonical public URL
        used for demos.
\end{itemize}

Trade-off: the cloudfront.net hostname is ugly. When a custom domain is
later added, the same distribution can be re-pointed to a Route 53 record
with an ACM cert; none of the downstream wiring changes.

\subsection{Changed: S3 bucket name is account-portable}

Bucket: \texttt{aikhub-audio-summaries-\$\{aws\_caller\_identity.account\_id\}}
(previously just \texttt{aikhub-audio-summaries}).

Rationale: S3 bucket names are a global namespace, and after a delete
AWS reserves the name for up to an hour. During the account migration
(see below) this blocked \texttt{terraform apply} for 19+ minutes until
the name was freed. Suffixing the bucket name with the account ID makes
every deploy unique and immune to this lock.

\subsection{Changed: Bedrock model upgraded to Claude Haiku 4.5}

\begin{itemize}[leftmargin=*]
  \item Was: \texttt{eu.anthropic.claude-3-haiku-20240307-v1:0}
        (flagged as \emph{Legacy} by Anthropic; returns
        \texttt{ResourceNotFoundException} on accounts that haven't used
        it in the last 30 days).
  \item Now: \texttt{eu.anthropic.claude-haiku-4-5-20251001-v1:0}.
  \item Prerequisite: the Anthropic \emph{use-case details} form must be
        submitted once per account through Bedrock $\to$ Model access
        $\to$ ``Submit use case details''. Applies to every Anthropic
        model; without it all Claude calls fail with
        ``Model use case details have not been submitted''.
\end{itemize}

\subsection{Changed: Audio storage uses object keys, not full URLs}

The \texttt{aikhub-papers} table now has an \texttt{audioKey} attribute
holding the S3 object key (e.g.\ \texttt{audio/2604.12911v1.mp3}) instead
of a full \texttt{audioUrl}. The email-sender Lambda reads
\texttt{S3\_BUCKET} from its environment and calls
\texttt{getSignedUrl(GetObjectCommand)} to mint a time-limited
pre-signed URL just before sending, with a TTL of seven days.

This replaces the previous scheme where emails linked to
\texttt{https://<bucket>.s3.amazonaws.com/...} URLs that 403'd because the
bucket blocks public access. The old ``S3 audio links are private''
entry under \textbf{Known Limitations} is now closed.

\subsection{Changed: Schema field renames}

\begin{itemize}[leftmargin=*]
  \item \texttt{aikhub-subscriptions.subscribedTopics} $\rightarrow$
        \texttt{topics}.
  \item Front-end \texttt{paper.paper\_id} $\rightarrow$
        \texttt{paper.externalId} (matches the DynamoDB PK; fixed
        \texttt{/papers/undefined} navigation).
  \item Backend \texttt{fetchPaperById(paper\_id)} $\rightarrow$
        \texttt{fetchPaperById(externalId)}.
\end{itemize}

One-shot migration scripts in \texttt{infra/scripts/} handled rewriting
the existing rows:
\texttt{add-paperids.mjs}, \texttt{backfill-schema.mjs},
\texttt{remove-old-paperid.mjs}.

\subsection{New: S3 Gateway VPC endpoint}

Added a second Gateway endpoint (\texttt{com.amazonaws.<region>.s3})
attached to the private route table. Private-subnet Lambdas now write
audio to S3 over the private backbone --- no NAT bandwidth, no egress
charges. DynamoDB and S3 are now both reached via free Gateway
endpoints.

\subsection{New: Frontend Dockerfile hardening}

The frontend image now runs \texttt{sed -i 's/\textbackslash r\$//'} on
\texttt{/docker-entrypoint.sh} before \texttt{chmod +x}. Without this
line, a checkout on Windows (CRLF line endings) produces a container
that crash-loops with \texttt{exec /docker-entrypoint.sh: no such file
or directory}. This was the root cause of a multi-hour debugging
session; the fix is one line and permanent.

\subsection{New: \texttt{empty\_bucket.py}}

A short boto3 script that deletes every object version + delete marker
in the audio bucket and then removes the bucket itself. Run before
\texttt{terraform destroy}: versioned buckets cannot be deleted while
they hold any non-current versions, and the AWS CLI one-liner for this
is fragile under Git Bash's JSON quoting rules.

\subsection{Incidents and recovery notes}

Captured here so future contributors don't spend the same evenings
debugging them.

\subsubsection*{Bedrock account-level quota = 0}

The initial AWS account had an \emph{applied} quota of zero tokens/day
for every Bedrock inference profile, with all ``Request increase'' paths
greyed out. Because the quota is account-level trust scoring (not a
per-model setting), no amount of model swapping fixes it. Remediation:
the stack was migrated to a teammate's account. The destroy-then-reapply
flow was rehearsed end-to-end and is reproducible from this document.

\subsubsection*{SES sandbox blocks arbitrary recipients}

Fresh accounts begin in SES sandbox, which requires \emph{both} sender
and recipient to be verified. For the demo, individual recipient email
addresses were verified through SES $\to$ Verified Identities. A
production-access request via Support Center $\to$ ``Service limit
increase'' $\to$ SES Sending Limits is the permanent fix; AWS typically
approves within 24 hours.

\subsubsection*{Wrong-account Terraform plan}

If \texttt{AWS\_PROFILE} is unset in a new shell, \texttt{terraform plan}
can read the default profile's credentials and show the entire stack as
\emph{deleted outside Terraform} --- tempting a disastrous re-apply into
the wrong account. Before any plan or apply, run
\texttt{aws sts get-caller-identity} and confirm the \texttt{Account}
field. Add \texttt{export AWS\_PROFILE=teammate} to
\texttt{\textasciitilde/.bashrc} to make the correct profile sticky.

\subsubsection*{Git Bash path mangling}

Commands that pass log-group names starting with \texttt{/}
(\texttt{/aws/lambda/aikhub-ai-processor}) are rewritten by MSYS into
Windows paths (\texttt{C:/Program Files/Git/aws/lambda/...}). Prefix
every such invocation with \texttt{MSYS\_NO\_PATHCONV=1}.

\subsubsection*{\texttt{charmap} codec error on Windows}

AWS CLI v2 on Windows fails to encode Unicode bytes returned by
\texttt{aws ec2 get-console-output}. Workaround: use boto3 with an
explicit \texttt{open(\ldots, encoding="utf-8")} when the output must
be persisted.

\subsection{Known Limitations --- updated}

Resolved since the initial document:
\begin{itemize}[leftmargin=*]
  \item \textbf{HTTPS}: now terminated at CloudFront.
  \item \textbf{Private audio links}: now pre-signed with a 7-day TTL at
        email-send time.
\end{itemize}

Still open:
\begin{itemize}[leftmargin=*]
  \item \textbf{Custom domain.} The public URL is still a
        \texttt{*.cloudfront.net} hostname. Trivial to add: buy the
        domain in Route 53, request an ACM cert in
        \texttt{us-east-1} (CloudFront requires us-east-1 certs), attach
        it to the distribution's \texttt{viewer\_certificate} block,
        point an A-alias at the distribution.
  \item \textbf{Unsubscribe flow.}
  \item \textbf{CI/CD and CloudWatch alarms.}
  \item \textbf{SES production access} (still in sandbox).
  \item \textbf{Bedrock IAM scope} --- \texttt{Resource = "*"} could be
        narrowed to the specific inference profile ARN.
\end{itemize}

\subsection{Updated cost snapshot}

The managed-compute tier adds a predictable floor on top of the
serverless costs documented earlier:

\begin{itemize}[leftmargin=*]
  \item \textbf{NAT Gateway} --- $\approx$\$35/month (unchanged; still
        the dominant fixed cost).
  \item \textbf{VPC interface endpoints} --- $\approx$\$16/month for SQS
        and SES. DynamoDB and S3 are Gateway endpoints and cost
        nothing.
  \item \textbf{ALB} --- $\approx$\$16/month base + per-LCU processing.
  \item \textbf{2$\times$ EC2 \texttt{t3.micro}} --- $\approx$\$15/month
        combined.
  \item \textbf{CloudFront} --- free up to 1\,TB/month of egress on the
        AWS Free Tier; after that \$0.085/GB. The distribution itself
        is free.
  \item \textbf{ECR} --- \$0.10/GB-month for stored images. With the
        keep-last-five lifecycle, $\approx$\$0.50/month.
\end{itemize}

Back-of-envelope idle cost with the full stack up:
\textbf{\$65--\$80/month}. Tear the stack down (\texttt{empty\_bucket.py}
then \texttt{terraform destroy}) to drop to $\approx$\$0 when not
demoing.

\section{Summary}

Every piece of back-end logic lives in one of four Lambdas, every piece
of state lives in DynamoDB or S3, and every interaction is mediated by
either SQS (for asynchronous work) or EventBridge (for scheduling).
The front-end and its Express API are containerised and served from an
auto-scaling EC2 fleet behind an Application Load Balancer, with public
HTTPS provided by a CloudFront distribution that fronts the ALB. The
entire stack --- serverless pipeline \emph{and} compute layer --- is
reproducible from \texttt{infra/} with \texttt{./build.sh \&\&
terraform apply}, and removable with \texttt{terraform destroy}
(preceded by \texttt{empty\_bucket.py}). The design is deliberately
boring: no custom infrastructure code, no cross-account trickery, no
hand-rolled retries --- just AWS primitives wired together so that each
component can be reasoned about, replaced, or scaled independently.

\end{document}
